% answers.tex —— 1977 年·江苏农垦系统初试卷 · 参考答案与详细解析片段
% 由 tools/generate_projects.py 依来源总账自 1977/1977_数学真题_江苏农垦系统初试_试题.md + 1977/1977_数学真题_江苏农垦系统初试_解析正文.md + 1977/1977_数学真题_江苏农垦系统初试_答案要点.md 生成，勿手改（改源 MD 后重跑）


\begin{gktip}
固定源 PDF《1977江苏农垦系统初试.pdf》为初试（预选）卷题面扫描件，未附解答，以下各题由编者独立推导。原卷题号为中文数字「一、」至「八、」，正文按题库规范作「第 N 题」；第 1 题含原卷自带的（1）（2）两小问，其余各题无小问。
\end{gktip}

\begin{answer}
\gksolhead{1}{}

\textbf{题目}：

计算：

（1） $\text{〔}1\frac{1}{24}-\left(\dfrac38+\dfrac16-\dfrac34\right)\times24\text{〕}\div(-5)^0\text{\text{；}}$

（2） $\lg\dfrac52+\lg\dfrac25-\lg10^3\text{.}$

\Needspace{4\baselineskip}
\textbf{解析}

\gksubpart{（1）}

\Needspace{5\baselineskip}
先算小括号内（通分为分母 $24$）：

\nopagebreak[4]
\[
\frac38+\frac16-\frac34=\frac{9}{24}+\frac{4}{24}-\frac{18}{24}=-\frac{5}{24}.
\]

\Needspace{5\baselineskip}
代入原式，注意 $(-5)^0=1$（底数非零）：

\nopagebreak[4]
\[
\text{〔}1\frac{1}{24}-\Bigl(-\frac{5}{24}\Bigr)\times24\text{〕}\div(-5)^0
=1\frac{1}{24}+5
=\frac{25}{24}+\frac{120}{24}
=\frac{145}{24}
=6\frac{1}{24}.
\]

\gksubpart{（2）}

\Needspace{5\baselineskip}
由对数运算法则，

\nopagebreak[4]
\[
\lg\frac52+\lg\frac25=\lg\Bigl(\frac52\times\frac25\Bigr)=\lg1=0,
\]

\Needspace{5\baselineskip}
故原式为

\nopagebreak[4]
\[
0-\lg10^3=-3.
\]

\gkanswer{答案（1）：$6\dfrac1{24}$（即 $\dfrac{145}{24}$）。}

\gkanswer{答案（2）：$-3$。}

\end{answer}
\begin{answer}
\gksolhead{2}{}

\textbf{题目}：

\Needspace{5\baselineskip}
解方程

\nopagebreak[4]
\[
\frac1{x+2}+\frac{4x}{x^2-4}-\frac2{x-2}=1.
\]

\Needspace{4\baselineskip}
\textbf{解析}

\Needspace{5\baselineskip}
分母要求 $x\ne\pm2$。两边同乘 $x^2-4$：

\nopagebreak[4]
\[
(x-2)+4x-2(x+2)=x^2-4,
\]

\Needspace{5\baselineskip}
整理得

\nopagebreak[4]
\[
3x-6=x^2-4,
\qquad
x^2-3x+2=0,
\]

\Needspace{5\baselineskip}
解得 $x=1$ 或 $x=2$。$x=2$ 使分母为零，为增根，舍去。代回验算：$x=1$ 时左端为

\nopagebreak[4]
\[
\frac13+\frac{4}{-3}-\frac{2}{-1}=\frac13-\frac43+2=1,
\]

与右端相等，故方程的解为 $x=1$。

\gkanswer{答案：$x=1$。}

\end{answer}
\begin{answer}
\gksolhead{3}{}

\textbf{题目}：

用长 $60$ 米的篱笆围成一矩形养鸡场地，一面利用原有旧墙，问矩形的长和宽各多少米才能使所围成的面积最大？

\Needspace{4\baselineskip}
\textbf{解析}

\Needspace{5\baselineskip}
设垂直于墙的一边（宽）为 $x$ 米（$0<x<30$），则平行于墙的一边（长）为 $(60-2x)$ 米，场地面积

\nopagebreak[4]
\[
S=x(60-2x)=-2(x-15)^2+450.
\]

当 $x=15$ 时 $S$ 取最大值 $450$，此时长为 $60-2\times15=30$（米）。即长 $30$ 米、宽 $15$ 米时所围面积最大，最大面积为 $450$ 平方米。

\gkanswer{答案：长 $30$ 米、宽 $15$ 米，最大面积 $450\ \mathrm{m^2}$。}

\end{answer}
\begin{answer}
\gksolhead{4}{}

\textbf{题目}：

把极坐标方程 $\rho=2\cos\theta$ 化成直角坐标方程，并说出它是什么曲线？

\Needspace{4\baselineskip}
\textbf{解析}

\Needspace{5\baselineskip}
在 $\rho=2\cos\theta$ 两边同乘 $\rho$，得 $\rho^2=2\rho\cos\theta$。由互化关系 $\rho^2=x^2+y^2$、$\rho\cos\theta=x$，得

\nopagebreak[4]
\[
x^2+y^2=2x,
\qquad\text{\text{即}}\qquad
(x-1)^2+y^2=1.
\]

所以原方程表示圆心在 $(1,0)$、半径为 $1$ 的圆（该圆过极点）。

\gkanswer{答案：$x^2+y^2-2x=0$，即 $(x-1)^2+y^2=1$；曲线是圆心在 $(1,0)$、半径为 $1$ 的圆。}

\end{answer}
\begin{answer}
\gksolhead{5}{}

\textbf{题目}：

在锐角 $\triangle ABC$ 中，已知 $\sin A=\dfrac35$，$\cos B=\dfrac5{13}$，求 $\cos C$ 的值。

\Needspace{4\baselineskip}
\textbf{解析}

题设 $\triangle ABC$ 为锐角三角形，故 $A$、$B$、$C$ 均为锐角。由 $\sin A=\dfrac35$ 得 $\cos A=\dfrac45$（锐角余弦取正值）；由 $\cos B=\dfrac5{13}$ 得 $\sin B=\dfrac{12}{13}$。

\Needspace{5\baselineskip}
由 $C=180^\circ-(A+B)$，

\nopagebreak[4]
\[
\cos C=-\cos(A+B)
=\sin A\sin B-\cos A\cos B
=\frac35\times\frac{12}{13}-\frac45\times\frac5{13}
=\frac{36-20}{65}
=\frac{16}{65}.
\]

$\cos C=\dfrac{16}{65}>0$，与 $C$ 为锐角相符。

\gkanswer{答案：$\cos C=\dfrac{16}{65}$。}

\end{answer}
\begin{answer}
\gksolhead{6}{}

\textbf{题目}：

三个数成等差数列，且它们的和是 $9$，如果把这个数列依次加上 $1$、$1$、$3$，那么所得的数又成等比数列。求这三个数。

\Needspace{4\baselineskip}
\textbf{解析}

\Needspace{5\baselineskip}
设这三个数为 $a-d$、$a$、$a+d$，由和为 $9$ 得 $a=3$。依次加上 $1$、$1$、$3$ 后得

\nopagebreak[4]
\[
4-d,\qquad 4,\qquad 6+d,
\]

\Needspace{5\baselineskip}
成等比数列，故中项的平方等于两端之积：

\nopagebreak[4]
\[
4^2=(4-d)(6+d)
\quad\Longrightarrow\quad
16=24-2d-d^2
\quad\Longrightarrow\quad
d^2+2d-8=0,
\]

解得 $d=2$ 或 $d=-4$。

$d=2$ 时，三个数为 $1,3,5$，加后得 $2,4,8$，公比 $2$；$d=-4$ 时，三个数为 $7,3,-1$，加后得 $8,4,2$，公比 $\dfrac12$。两组各项均非零且确成等比，均符合题意。

\gkanswer{答案：这三个数为 $1,3,5$ 或 $7,3,-1$。}

\end{answer}
\begin{answer}
\gksolhead{7}{}

\textbf{题目}：

求以椭圆 $\dfrac{x^2}{49}+\dfrac{y^2}{25}=1$ 的焦点和长轴的端点作为顶点和焦点的双曲线方程。

\Needspace{4\baselineskip}
\textbf{解析}

\Needspace{5\baselineskip}
椭圆 $\dfrac{x^2}{49}+\dfrac{y^2}{25}=1$ 中，$a^2=49$、$b^2=25$，故

\nopagebreak[4]
\[
c^2=a^2-b^2=49-25=24,
\]

焦点为 $(\pm2\sqrt6,0)$，长轴端点为 $(\pm7,0)$。

\Needspace{5\baselineskip}
所求双曲线以椭圆的焦点为顶点、长轴端点为焦点，焦点仍在 $x$ 轴上，故

\nopagebreak[4]
\[
a'^2=24,
\qquad
c'^2=49,
\qquad
b'^2=c'^2-a'^2=49-24=25.
\]

\Needspace{5\baselineskip}
于是双曲线方程为

\nopagebreak[4]
\[
\frac{x^2}{24}-\frac{y^2}{25}=1.
\]

\gkanswer{答案：$\dfrac{x^2}{24}-\dfrac{y^2}{25}=1$。}

\end{answer}
\begin{answer}
\gksolhead{8}{}

\textbf{题目}：

等腰三角形 $ABC$ 的顶角 $A$ 为 $100^\circ$，底角 $B$ 的平分线交腰 $AC$ 于 $D$，求证：$BC=AD+BD$。

\Needspace{4\baselineskip}
\textbf{解析}

\Needspace{5\baselineskip}
设腰长 $AB=AC=a$。顶角 $A=100^\circ$，故底角

\nopagebreak[4]
\[
\angle B=\angle C=\frac{180^\circ-100^\circ}{2}=40^\circ.
\]

$BD$ 平分 $\angle B$，所以 $\angle ABD=\angle DBC=20^\circ$；在 $\triangle ABD$ 中 $\angle ADB=180^\circ-100^\circ-20^\circ=60^\circ$。

\Needspace{5\baselineskip}
在 $\triangle ABD$ 中用正弦定理（$AD$ 对 $20^\circ$、$BD$ 对 $100^\circ$、$AB$ 对 $60^\circ$）：

\nopagebreak[4]
\[
\frac{AD}{\sin20^\circ}=\frac{BD}{\sin100^\circ}=\frac{a}{\sin60^\circ}.
\]

\Needspace{5\baselineskip}
在 $\triangle ABC$ 中用正弦定理（$BC$ 对 $100^\circ$、$AB$ 对 $40^\circ$）：

\nopagebreak[4]
\[
\frac{BC}{\sin100^\circ}=\frac{a}{\sin40^\circ}.
\]

\Needspace{5\baselineskip}
于是（和差化积：$\sin20^\circ+\sin100^\circ=2\sin60^\circ\cos40^\circ$）

\nopagebreak[4]
\[
AD+BD
=\frac{a(\sin20^\circ+\sin100^\circ)}{\sin60^\circ}
=\frac{a\cdot2\sin60^\circ\cos40^\circ}{\sin60^\circ}
=2a\cos40^\circ,
\]

\Needspace{5\baselineskip}
而

\nopagebreak[4]
\[
BC=\frac{a\sin100^\circ}{\sin40^\circ}
=\frac{a\sin80^\circ}{\sin40^\circ}
=\frac{a\cdot2\sin40^\circ\cos40^\circ}{\sin40^\circ}
=2a\cos40^\circ.
\]

\Needspace{5\baselineskip}
两式相等，故

\nopagebreak[4]
\[
BC=AD+BD.
\]

\gkanswer{答案：见解析版}

\end{answer}